\chapter{Repo and Specials}\label{ch:m2:repo-and-specials}

On Monday 16 September 2019 American companies paid their quarterly taxes
and USD~54 billion of newly auctioned Treasury notes and bonds settled. Both
payments drained cash from banks' accounts at the Federal Reserve: reserves fell
by about 120 billion dollars in two days, to 1.34 trillion, their lowest since
2012. On Tuesday the rate on overnight loans secured by Treasuries, the market
where every dealer finances its inventory, printed at 5.25\%. The Federal
Reserve's target range for its policy rate was 2.00 to 2.25\%, and the federal
funds rate itself rose above it, to 2.30\%. At 9:30 that morning the New York
Fed offered up to 75 billion dollars in an overnight repo operation and lent
53 billion. Repo is
the plumbing under every position in this book: a bond bought is a bond
financed, and a bond sold short is a bond borrowed. This chapter follows cash
and collateral through it, and prices the one bond that everyone wants to
borrow.

\section{The repo trade from both sides}

A repurchase agreement (One Quant Book 1, chapter 6) is a sale of securities
with an agreement to buy them back at a fixed price on a later date: an
exchange of collateral for cash, reversed later, whose price difference is
interest.

\begin{definition}[Repo rate, reverse repo, term repo]\label{def:m2:repo-and-specials:rate}
The \emph{repo rate}\index{repo rate} is the simple interest rate, actual/360
in dollars, implied by the difference between the repurchase and the sale
price. The side that sells the collateral and borrows cash does a repo; the
side that lends cash and receives collateral does a \emph{reverse
repo}\index{reverse repo}. An overnight repo matures on the next business day;
a \emph{term repo}\index{term repo} has a fixed longer maturity; an open repo
rolls each day until one side ends it.
\end{definition}

\begin{example}[Financing the ten-year note]\label{ex:m2:repo-and-specials:finance}
A dealer buys USD~100 million of the ten-year note of
\cref{ch:m2:government-bonds} at a dirty price of 100.870911 and finances it
in repo for 30 days at 3.90\%. With an illustrative 2\% haircut, the lender
pays $100\,000\,000 \times 1.00870911 \times 0.98 = \text{USD}~98\,853\,492$ and
receives 321\,274 of interest at the end; the dealer funds the other 2\% from
its own capital. Over the 30 days the note accrues USD~346\,467 of coupon.
\end{example}

\begin{proposition}[Carry of a financed bond]\label{prop:m2:repo-and-specials:carry}
A long position of face $N$ in a bond with coupon $c$ and dirty price $P$,
financed in full at the \omterm{def:m2:repo-and-specials:rate}{repo rate} $r$ for $d$ days in a coupon period of $D$
days, earns, if yields do not move,
\[
\text{carry} \;=\; N\,\frac{c}{f}\,\frac{d}{D} \;-\; N\,\frac{P}{100}\,r\,\frac{d}{360},
\]
plus the (small) pull of the clean price towards par. The yield rise that
would wipe it out is the carry divided by the position's DV01.
\end{proposition}
\begin{proof}
The first term is the coupon accrued over the $d$ days, which the holder
collects in the dirty price; the second is the repo interest on the cash
borrowed. A yield change $\Delta y$ changes the value by $-\mathrm{DV01}\times
\Delta y$ in basis points.
\end{proof}

At 3.90\% the dealer's carry on the ten-year is $346\,467 - 327\,830 =
\text{USD}~18\,637$ over the month, and a rise of 0.23 basis points in yield would
erase it. Positive carry is the reward for owning a bond that yields more
than its financing; a curve on which long bonds yield less than overnight
money makes carry negative, and holding bonds a bet on falling yields.

\section{Tri-party, bilateral and sponsored}

\begin{definition}[Tri-party and bilateral repo]\label{def:m2:repo-and-specials:triparty}
In \emph{tri-party repo}\index{tri-party repo} a clearing bank acting as
agent holds the collateral, allocates it from the borrower's inventory
according to the lender's eligibility schedule, values it and applies the
haircuts; the lender is lent against a basket, never a specific security.
In \emph{bilateral repo}\index{bilateral repo} the two sides settle directly,
delivery versus payment, and the cash lender can ask for a specific security.
\end{definition}

\begin{definition}[Sponsored repo]\label{def:m2:repo-and-specials:sponsored}
In \emph{sponsored repo}\index{sponsored repo} a dealer that is a member of
the Treasury clearing house sponsors a non-member, typically a \omterm{def:m2:money-markets:mmf}{money-market
fund} lending cash or a hedge fund borrowing it, into the clearing house: the
trade is novated to the central counterparty and nets against the dealer's
other cleared trades.
\end{definition}

\begin{omfigure}
\begin{tikzpicture}[font=\footnotesize,
  box/.style={draw,thick,rounded corners=2pt,align=center,minimum height=0.8cm,minimum width=1.8cm,inner sep=3pt},
  arr/.style={-{Stealth[length=2mm]},thick}]
% bilateral
\node[font=\small\bfseries] at (1.6,2.45) {bilateral};
\node[box,draw=omDef,fill=omDef!8] (l1) at (0.3,1.2) {cash lender};
\node[box,draw=omThm,fill=omThm!8] (b1) at (2.9,1.2) {dealer};
\draw[arr,omDef] (l1.20) -- node[above=2pt] {cash} (b1.160);
\draw[arr,omThm] (b1.200) -- node[below=2pt,align=center] {specific\\security} (l1.340);
% tri-party
\node[font=\small\bfseries] at (6.6,2.45) {tri-party};
\node[box,draw=omDef,fill=omDef!8] (l2) at (5.3,1.2) {cash lender};
\node[box,draw=omThm,fill=omThm!8] (b2) at (7.9,1.2) {dealer};
\node[box,draw=omProp,fill=omProp!8] (a2) at (6.6,-0.6) {agent bank:\\allocates, values};
\draw[arr,omDef] (l2) -- node[above=9pt] {cash} (b2);
\draw[arr,omThm] (b2) -- node[right=2pt] {basket} (a2);
\draw[arr,omProp] (a2) -- node[left=2pt,align=center] {eligible\\collateral} (l2);
% sponsored
\node[font=\small\bfseries] at (11.6,2.45) {sponsored};
\node[box,draw=omDef,fill=omDef!8] (l3) at (10.3,1.2) {money fund};
\node[box,draw=omThm,fill=omThm!8] (b3) at (12.9,1.2) {hedge fund};
\node[box,draw=omProp,fill=omProp!8] (c3) at (11.6,-0.6) {clearing house\\(dealer sponsors both)};
\draw[arr,<->,omDef] (l3) -- node[left=2pt] {novated} (c3);
\draw[arr,<->,omThm] (b3) -- node[right=2pt] {novated} (c3);
\end{tikzpicture}

\omcaption{Three ways to do the same repo. Bilaterally, the cash lender can
ask for a specific security. In tri-party, an agent bank allocates collateral
from the dealer's inventory to fit the lender's schedule. Sponsored, a
dealer brings both sides into the clearing house, which faces each of them,
and the dealer's two trades net.}\label{fig:m2:repo-and-specials:structures}
\end{omfigure}

Netting is the point. A dealer that borrows cash from a money fund in one repo
and lends it to a hedge fund in another holds two trades on its balance sheet
if both are bilateral, and, if both are cleared through the same central
counterparty, a net position close to zero. Because a dealer's balance sheet
is costly, above all on the dates it is measured (\cref{ch:m2:money-markets}),
sponsored clearing has grown into a major way for dealers to intermediate
repo at scale (\cref{fig:m2:repo-and-specials:sponsored}), and the clearing
mandate of \cref{dat:m2:the-treasury-market:clearing} will send much more
repo through the clearing house.

\begin{omfigure}
\begin{tikzpicture}
\begin{axis}[width=11cm,height=4.8cm,
  ylabel={USD trillion},
  tick label style={font=\small},label style={font=\small},
  xmin=2020,xmax=2027,ymin=0,ymax=1.8,grid=major,grid style={gray!25},
  xtick={2020,2021,2022,2023,2024,2025,2026,2027},
  xticklabel style={/pgf/number format/1000 sep=},
  legend columns=2,legend style={font=\footnotesize,at={(0.5,-0.22)},anchor=north,draw=none,fill=none,/tikz/every even column/.append style={column sep=8pt}},
  table/col sep=comma]
\addplot[omDef,very thick] table[x=t,y=reverse]{figdata/markets-2/05-repo-and-specials/sponsored.csv};
\addplot[omThm,very thick] table[x=t,y=repo]{figdata/markets-2/05-repo-and-specials/sponsored.csv};
\legend{sponsored reverse repo (cash lent by sponsored members),sponsored repo (cash borrowed)}
\end{axis}
\end{tikzpicture}

\omcaption{FICC \omterm{def:m2:repo-and-specials:sponsored}{sponsored repo} volumes at each month-end, March 2020 to August
2026. Money funds lend through sponsored \omterm{def:m2:repo-and-specials:rate}{reverse repo}, hedge funds borrow
through \omterm{def:m2:repo-and-specials:sponsored}{sponsored repo}; both grew four- to fivefold, with peaks at year-ends,
when a netted trade is worth most to a dealer. Data: Office of Financial
Research, Hedge Fund Monitor (FICC data).}\label{fig:m2:repo-and-specials:sponsored}
\end{omfigure}

\begin{dated}{2026-09}{FICC sponsored repo}\label{dat:m2:repo-and-specials:sponsored}
\omterm{def:m2:repo-and-specials:sponsored}{Sponsored repo} plus sponsored \omterm{def:m2:repo-and-specials:rate}{reverse repo} reached USD~2.96 trillion on 31
December 2025 (1.38 trillion of repo and 1.58 trillion of \omterm{def:m2:repo-and-specials:rate}{reverse repo}), and
stood at 2.33 trillion on 21 August 2026, the latest date published.
\end{dated}

\section{Specials and specialness}

Most repo is \emph{general collateral} (One Quant Book 1, chapter 16): the cash
lender accepts any security of a class, and the rate is the price of cash.
When someone needs a particular security, to deliver it against a short sale
or into a futures contract, the rate is also the price of that security.

\begin{definition}[Special repo rate and specialness]\label{def:m2:repo-and-specials:special}
The \emph{special repo rate}\index{special repo rate} of a security is the
\omterm{def:m2:repo-and-specials:rate}{repo rate} at which cash lenders will lend against that specific security. Its
\emph{specialness}\index{specialness} is the general-collateral rate minus the
special rate: what a cash lender gives up in interest to obtain the security.
\end{definition}

\omterm{def:m2:repo-and-specials:special}{Specialness} is paid to the holder of the security: a dealer who owns the
on-the-run ten-year can borrow cash against it more cheaply than against any
other note. The on-the-run is the most shorted issue, since traders hedge with
it, and the most often special. As its successor's auction approaches, the
shorts move to the new issue and its \omterm{def:m2:repo-and-specials:special}{specialness} fades
(\cref{fig:m2:repo-and-specials:special}).

\begin{proposition}[What specialness is worth]\label{prop:m2:repo-and-specials:value}
A holder who finances at \omterm{def:m2:repo-and-specials:special}{specialness} $s_k$ for $d_k$ days in successive
periods saves, per 100 of dirty price $P$,
\[
V \;=\; P\,\sum_k \frac{d_k\,s_k}{360}
\]
price points. The security can therefore trade that much richer than an
otherwise identical one, or $V/\mathrm{DV01}$ basis points lower in yield,
before owning it stops paying.
\end{proposition}
\begin{proof}
Each day the holder repos the security at $\mathrm{GC} - s_k$ instead of GC,
saving $P\,s_k/360$ per 100. Summed over the expected special life this is the
present gain, to first order; an investor indifferent between the two
securities pays it up front in the price.
\end{proof}

\begin{omfigure}
\begin{tikzpicture}
\begin{axis}[width=11cm,height=4.5cm,
  xlabel={days after issue},ylabel={specialness (bp)},
  tick label style={font=\small},label style={font=\small},
  xmin=0,xmax=90,ymin=0,ymax=55,grid=major,grid style={gray!25},
  table/col sep=comma]
\addplot[omDef,very thick,const plot] table[x=day,y=bp]{figdata/markets-2/05-repo-and-specials/special.csv};
\node[font=\footnotesize,anchor=south west] at (axis cs:2,46) {new issue, heavily shorted};
\node[font=\footnotesize,anchor=south east] at (axis cs:88,11) {next auction approaches};
\end{axis}
\end{tikzpicture}

\omcaption{An illustrative \omterm{def:m2:repo-and-specials:special}{specialness} path for a new ten-year note over the
three months to its successor's auction: 45 basis points in the first month,
25 in the second, 10 in the third. The area under the step is the value of
\cref{prop:m2:repo-and-specials:value}. Levels are illustrative. Data: the
chapter's weekend problem.}\label{fig:m2:repo-and-specials:special}
\end{omfigure}

\section{Fails and the fails charge}

A short who cannot borrow a security it has sold fails to deliver it. Until
2009 the convention was simply to postpone delivery, without penalty and at
an unchanged price: failing cost only the interest on the sale proceeds, which
the buyer does not pay until delivery. When short rates collapsed after
Lehman Brothers' insolvency in September 2008, that cost collapsed with them
and fails reached an extraordinary volume. The market's practice committee
then introduced a charge, from May 2009.

\begin{definition}[Fails charge]\label{def:m2:repo-and-specials:fails}
The \emph{fails charge}\index{fails charge} on a Treasury transaction that
fails to settle is, for each day of the fail,
\[
C \;=\; \frac{1}{360}\times 0.01 \times \max(3 - R,\,0) \times P,
\]
where $P$ is the proceeds and $R$, in percent, is the reference rate (the lower
limit of the Federal Reserve's target range) on the preceding business day. It
is paid by the failing seller to the buyer.
\end{definition}

\begin{proposition}[A floor under special rates]\label{prop:m2:repo-and-specials:floor}
With a \omterm{def:m2:repo-and-specials:fails}{fails charge} of $(3\% - R)^+$ a year, no short borrows a security at a
special rate below $-(3\% - R)^+$: rather than accept less, it fails. When
$R \ge 3\%$ the floor is zero; when $R < 3\%$ special rates can be negative,
down to $R - 3\%$.
\end{proposition}
\begin{proof}
To deliver, the short lends cash $C$ at the special rate $s$ and receives the
security, then delivers it and is paid about $C$: it earns $s$ on $C$. Failing
instead, it is not paid, so earns nothing, and pays $(3\% - R)^+$ on $C$. It
borrows the security if and only if $s > -(3\% - R)^+$.
\end{proof}

\begin{dated}{2026-09}{The fails charge today}\label{dat:m2:repo-and-specials:fails}
The charge applies to Treasury transactions entered into since 1 May 2009.
With the lower limit of the target range at 3.75\% since 17 September 2026,
$\max(3 - 3.75, 0) = 0$: failing to deliver a Treasury currently carries no
charge, and special rates are floored at zero by the proposition. The charge
bites again if the range's lower limit falls below 3\%.
\end{dated}

\section{September 2019}

Reserves had been falling since the Federal Reserve began to shrink its
balance sheet, and kept falling after that programme ended in August 2019;
at the same time the Treasury's growing issuance had filled dealers'
inventories with securities to finance. On 16 September the tax date and the coupon settlement took
another 120 billion of reserves out of the system in two days. The overnight
Treasury \omterm{def:m2:repo-and-specials:rate}{repo rate} reached 5.25\% on 17 September
(\cref{fig:m2:repo-and-specials:sept2019}), and the federal funds rate left
its target range. The New York Fed's repo operation at 9:30 that morning
added 53 billion of reserves; the Committee cut its range on the 18th, and
from mid-October the Federal Reserve bought bills, at about 60 billion dollars
a month at first, to keep reserves at least at their early-September level.

\begin{omfigure}
\begin{tikzpicture}
\begin{axis}[width=11cm,height=5cm,
  xlabel={September--October 2019},ylabel={rate (\%)},
  tick label style={font=\small},label style={font=\small},
  xmin=0,xmax=42,ymin=1.5,ymax=5.6,grid=major,grid style={gray!25},
  xtick={2,9,16,17,23,30,37},xticklabels={2 Sep,9,16,,23,30 Sep,7 Oct},
  legend columns=3,legend style={font=\footnotesize,at={(0.5,-0.3)},anchor=north,draw=none,fill=none,/tikz/every even column/.append style={column sep=8pt}},
  table/col sep=comma]
\addplot[omThm,very thick,mark=*,mark size=1.2pt] table[x=t,y=sofr]{figdata/markets-2/05-repo-and-specials/sept2019.csv};
\addplot[omDef,very thick] table[x=t,y=effr]{figdata/markets-2/05-repo-and-specials/sept2019.csv};
\addplot[gray,thick,dashed,const plot] table[x=t,y=upper]{figdata/markets-2/05-repo-and-specials/sept2019.csv};
\addplot[gray,thick,dashed,const plot,forget plot] table[x=t,y=lower]{figdata/markets-2/05-repo-and-specials/sept2019.csv};
\legend{SOFR (Treasury repo),effective federal funds rate,target range}
\node[font=\footnotesize,anchor=west] at (axis cs:18,4.9) {17 September: 5.25\%};
\end{axis}
\end{tikzpicture}

\omcaption{Overnight rates around 17 September 2019. The secured \omterm{def:m2:repo-and-specials:rate}{repo rate}
jumped far above the policy range; the federal funds rate, an unsecured bank
market, broke its ceiling by five basis points. After the Federal Reserve's
operations and its cut of the 18th the rates fell back inside the new range,
except for a quarter-end spike on 30 September. Data: Federal Reserve Bank of
New York and Board of Governors, via FRED.}\label{fig:m2:repo-and-specials:sept2019}
\end{omfigure}

The episode showed that a \omterm{def:m2:central-banks-and-the-short-rate:corridor}{floor system} (\cref{ch:m2:central-banks-and-the-short-rate})
needs its reserves to be ample in every bank that has to lend them, not only on
average. In July 2021 the Committee created a standing repo facility, a daily
overnight repo operation against Treasuries at a fixed minimum rate, which now
puts a ceiling on the \omterm{def:m2:repo-and-specials:rate}{repo rate} for eligible counterparties.

\section{Tutorial: a repo book}\label{tut:m2:repo-and-specials:book}

\begin{tutorial}
\textbf{Goal.} Book the financing of \cref{ex:m2:repo-and-specials:finance},
compute its carry and its margin call after a price fall, and value the
\omterm{def:m2:repo-and-specials:special}{specialness} path of \cref{fig:m2:repo-and-specials:special}. \textbf{End
state:} the numbers of the example and of the weekend problem.
\begin{tutsteps}
\item \textbf{The trade}: cash is collateral value less the haircut; interest
accrues actual/360.
\omcode{code/firm/repo/firm_repo.py}{11}{32}{A repo trade: cash, interest, repurchase price.}\label{lst:m2:repo-and-specials:trade}
\item \textbf{Margin.} If the collateral loses value, the cash borrower must
post more to restore the haircut.
\omcode{code/firm/repo/firm_repo.py}{39}{45}{The margin call on a repo after a price move.}\label{lst:m2:repo-and-specials:margin}
A fall of 1.5 points on USD~100 million calls USD~1.5 million on the day.
\item \textbf{\omterm{def:m2:repo-and-specials:special}{Specialness} and fails.}
\omcode{code/firm/repo/firm_repo.py}{52}{65}{The value of specialness, the fails charge, and the floor it puts under special rates.}\label{lst:m2:repo-and-specials:special}
Check the practice committee's example: USD~5\,555.56 a day on 100 million
when the reference rate is 1\%.
\item \textbf{Run} \texttt{repo\_demo.financed\_note()} and
\texttt{repo\_demo.special\_value()}; \texttt{fig\_repo.py} writes the three
charts.
\end{tutsteps}
\textbf{What to change next.} Finance the note at 4.30\% and find the
negative carry; then let the \omterm{def:m2:repo-and-specials:special}{specialness} of the first month double and
express the new value in 32nds.
\end{tutorial}

\section{Build: the repo book}\label{bld:m2:repo-and-specials:book}

\begin{build}
\bfield{Purpose} Every position the miniature firm holds in bonds is
financed or borrowed; the repo book records the trades, accrues them, calls
and meets margin, and tells the traders what their bonds cost to hold and what
their shorts cost to borrow.
\bfield{Interface} \texttt{RepoTrade(side, collateral, face, dirty\_price,
haircut, rate, start, end)} with \texttt{cash}, \texttt{interest(until)},
\texttt{repurchase\_price()}; \texttt{margin\_call(trade, new\_dirty\_price,
accrued\_to)}; \texttt{specialness(gc, special)};
\texttt{value\_of\_specialness(dirty, path)}; \texttt{fails\_charge(proceeds,
R, days)}; \texttt{special\_rate\_floor(R)}; \texttt{carry(\ldots)}.
\bfield{Rules} Actual/360; the haircut withheld from the cash; margin restores
the haircut on cash plus accrued interest; the \omterm{def:m2:repo-and-specials:fails}{fails charge} by the practice
committee's formula with the reference rate of the preceding business day.
\bfield{Acceptance tests} \texttt{code/firm/repo/tests/}: cash, interest and
repurchase price; margin calls in both directions; the practice committee's
worked example to the cent; the special-rate floor above and below 3\%; a bond
whose coupon equals its financing has zero carry.
\bfield{Stretch} Open repos repriced daily from a rate feed; substitution of
collateral; a tri-party eligibility schedule; interest accrual through the
overnight-rate engine (\cref{bld:m2:central-banks-and-the-short-rate:rfr}).
\end{build}

\begin{omsources}
\item Board of Governors of the Federal Reserve System, FEDS Notes, ``What
Happened in Money Markets in September 2019?'', 27 February 2020.
\item Treasury Market Practices Group, \emph{Frequently Asked Questions: TMPG
Fails Charges}, 2013.
\item Office of Financial Research, \emph{Hedge Fund Monitor}: FICC sponsored
repo service volumes.
\item K.\,D.\ Garbade, F.\,M.\ Keane, L.\ Logan, A.\ Stokes and J.\ Wolgemuth,
``The introduction of the TMPG fails charge for U.S. Treasury securities'',
Federal Reserve Bank of New York \emph{Economic Policy Review}, 2010.
\end{omsources}

\section{Exercises}

\begin{exercise}[$\star$]\label{exo:m2:repo-and-specials:1}
A fund lends USD~50 million in repo from Friday to Tuesday at 3.90\%. What
interest does it receive?
\end{exercise}

\begin{exercise}[$\star$]\label{exo:m2:repo-and-specials:2}
USD~100 million face of a note with a dirty price of 101.5 is repoed with a 2\%
haircut. How much cash does the borrower receive?
\end{exercise}

\begin{exercise}[$\star$]\label{exo:m2:repo-and-specials:3}
General collateral trades at 3.90\% and a note's special rate at 3.10\%. Give
its \omterm{def:m2:repo-and-specials:special}{specialness}, and what owning USD~100 million of it saves in financing per
day.
\end{exercise}

\begin{exercise}[$\star\star$]\label{exo:m2:repo-and-specials:4}
Verify the carry of the dealer's ten-year position in
\cref{ex:m2:repo-and-specials:finance} and the yield rise that would erase it.
What is the carry if the \omterm{def:m2:repo-and-specials:rate}{repo rate} is 4.30\%?
\end{exercise}

\begin{exercise}[$\star\star$]\label{exo:m2:repo-and-specials:5}
A seller fails to deliver USD~100 million of a note for five days. What does it
owe under the \omterm{def:m2:repo-and-specials:fails}{fails charge} when the reference rate is 1\%? When it is 3.75\%?
\end{exercise}

\begin{exercise}[$\star\star$]\label{exo:m2:repo-and-specials:6}
In 2021 the reference rate was 0\%. What was the lowest special rate a short
would accept rather than fail? What is it today? What happened to specials
before 2009 when rates were near zero?
\end{exercise}

\begin{exercise}[$\star\star\star$]\label{exo:m2:repo-and-specials:7}
\emph{Coding.} With \texttt{margin\_call}, compute the call on the dealer's
30-day repo of \cref{ex:m2:repo-and-specials:finance} if the note's dirty price
falls by 1.5 points on the first day. Then compute the call after 30 days if
the price has not moved, counting the accrued repo interest.
\end{exercise}

\begin{exercise}[$\star\star\star$]\label{exo:m2:repo-and-specials:8}
\emph{Find the flaw.} ``The federal funds rate traded above the top of the
range on 17 September 2019, so the Committee had tightened policy.'' Correct
the statement from \cref{fig:m2:repo-and-specials:sept2019} and the chapter.
\end{exercise}

\section{Problem: The Special}

\begin{problem}[{Weekend problem --- what a special is worth}]\label{pb:m2:repo-and-specials:1}
The ten-year note of \cref{ex:m2:repo-and-specials:finance} is the
on-the-run. A dealer expects it to trade special by 45 basis points for the
next 30 days, 25 for the 30 after, and 10 for the last 30 before the next
auction, with general collateral at 3.90\%. Its dirty price is 100.870911 and its
DV01 is 0.080444 per 100.

\medskip\noindent\textbf{Part I --- The rates.}
\begin{enumerate}
\item Give the \omterm{def:m2:repo-and-specials:special}{special repo rate} in each of the three months.
\item What does the dealer save in financing on USD~100 million in the first
month?
\item And over the 90 days?
\item Why is the \omterm{def:m2:repo-and-specials:special}{specialness} highest just after issue?
\item Who pays the \omterm{def:m2:repo-and-specials:special}{specialness}, and why would they?
\end{enumerate}

\medskip\noindent\textbf{Part II --- The price.}
\begin{enumerate}[resume]
\item Give the value of the \omterm{def:m2:repo-and-specials:special}{specialness} in price points per 100
(\cref{prop:m2:repo-and-specials:value}).
\item Express it in 32nds.
\item Express it in basis points of yield.
\item An off-the-run note of nearly the same maturity and coupon yields 4.210\%.
Ignoring liquidity, at what yield should the on-the-run trade?
\item It trades at 4.195\%. Is it rich or cheap on financing alone?
\end{enumerate}

\medskip\noindent\textbf{Part III --- The trade.}
\begin{enumerate}[resume]
\item Describe the trade that captures the difference, leg by leg, including
how each leg is financed.
\item What does each leg earn or pay in repo?
\item What if \omterm{def:m2:repo-and-specials:special}{specialness} turns out to be twice the dealer's forecast?
\item What if a large short squeezes the issue and its special rate goes to
zero?
\item With the reference rate at 3.75\%, how low can the special rate go?
\end{enumerate}

\medskip\noindent\textbf{Part IV --- Judgement.}
\begin{enumerate}[resume]
\item Why is the on-the-run's yield discount larger than its financing value
alone would justify, in most markets?
\item Why does a dealer that holds the on-the-run and lends it out care about
the \omterm{def:m2:repo-and-specials:fails}{fails charge}?
\item How would the Treasury's decision to reopen the issue (sell more of it)
change the \omterm{def:m2:repo-and-specials:special}{specialness}?
\item State the \emph{named result}: the value of the expected \omterm{def:m2:repo-and-specials:special}{specialness} in
price, 32nds and yield.
\item In one sentence: what is a special?
\end{enumerate}
\end{problem}

\section{Interview questions}

\begin{interviewq}[$\star$ \iqroles{trader, bank}]\label{iq:m2:repo-and-specials:1}
What is repo, and why is it the most important funding market for bond
traders?
\end{interviewq}

\begin{interviewq}[$\star$ \iqroles{trader, researcher}]\label{iq:m2:repo-and-specials:2}
What does it mean for a bond to trade special, and who benefits?
\end{interviewq}

\begin{interviewq}[$\star\star$ \iqroles{trader, researcher, bank}]\label{iq:m2:repo-and-specials:3}
What happened in the repo market in September 2019, and what did the Fed do?
\end{interviewq}

\begin{interviewq}[$\star\star$ \iqroles{trader}]\label{iq:m2:repo-and-specials:4}
You are long a ten-year note financed in repo. What is your carry, and what is
your breakeven?
\end{interviewq}

\begin{interviewq}[$\star\star$ \iqroles{bank, developer}]\label{iq:m2:repo-and-specials:5}
Why do dealers prefer sponsored, cleared repo to \omterm{def:m2:repo-and-specials:triparty}{bilateral repo}?
\end{interviewq}

\begin{interviewq}[$\star\star\star$ \iqroles{researcher, trader}]\label{iq:m2:repo-and-specials:6}
Can a \omterm{def:m2:repo-and-specials:rate}{repo rate} be negative? Derive the lower bound on a special rate.
\end{interviewq}
